% TOP_PROBLEM: {"id":30007681,"problem_number":"LOCAL-30007681","title":"Exclusionary contracting","category_id":21,"category":{"id":21,"name":"economics","display_name":"Economics","description":"Open economic theory statements, published research agendas, and proposed scoped specifications; question provenance and historical priority are qualified per record.","slug":"economics","order_index":21,"created_at":"2026-10-08T00:00:00Z"},"status":"research_question","external_url":null,"legacy_ids":[],"published":false,"statement":"Compare exclusive-contract regimes in a specified supply-chain model, allowing coordination, entry deterrence and supplier rematching.","economics":{"original_id":170,"field":"Industrial organization and competition","kind":"theoretical","formalization_basis":"adapted_research_question","source_status":"research_agenda","priority_status":"earliest_found","priority_note":"Earliest relevant explicit question or agenda passage located in this audit. No claim of first-ever formulation; earlier versions and later solutions were not exhaustively traced.","earliest_verified_reference":{"authors":["Robin S. Lee","Michael D. Whinston","Ali Yurukoglu"],"title":"Structural Empirical Analysis of Contracting in Vertical Markets","year":2021,"url":"https://robinlee.sites.fas.harvard.edu/papers/HandbookIO_Vol4_VerticalMarkets.pdf","doi":"10.3386/w29282","locator":"Section 5, printed pp. 58–59, future research directions","evidence_summary":"Richer contracts, partner choice and investment remain active modeling directions."},"current_reference":{"authors":["Robin S. Lee","Michael D. Whinston","Ali Yurukoglu"],"title":"Structural Empirical Analysis of Contracting in Vertical Markets","year":2021,"url":"https://robinlee.sites.fas.harvard.edu/papers/HandbookIO_Vol4_VerticalMarkets.pdf","doi":"10.3386/w29282","locator":"Section 5, printed pp. 58–59, future research directions","evidence_summary":"Richer contracts, partner choice and investment remain active modeling directions."},"formalization_note":"The source verifies a broader question or research agenda; the population, estimand, equations and restrictions here are our scoped adaptation. This is not a quoted canonical conjecture, and the audit does not prove contemporary unsolved status for every special case. Independent math review tightened feasibility, existence, or estimand assumptions where required."},"statusQualification":"Published question or research agenda verified at the cited locator. The mathematical specification is adapted for this catalogue and includes additional assumptions; source publication does not establish first-ever priority or certify this exact model remains unsolved. The source verifies a broader question or research agenda; the population, estimand, equations and restrictions here are our scoped adaptation. This is not a quoted canonical conjecture, and the audit does not prove contemporary unsolved status for every special case. Independent math review tightened feasibility, existence, or estimand assumptions where required.","statusReviewedAt":"2026-10-08"}
% FIRST_POSTED: 2026-10-08
% ENTRY_KIND: statement_only
% !TeX program = lualatex
\documentclass[11pt]{article}
\usepackage{fontspec}
\usepackage[a4paper,margin=25mm]{geometry}
\usepackage{amsmath,amssymb,mathtools}
\usepackage{xurl}
\usepackage[hidelinks]{hyperref}
\newcommand{\sref}[2]{\hyperlink{source:#1:#2}{[S#2]}}
\setlength{\emergencystretch}{3em}
\begin{document}
% problemId:30007681
\section*{Exclusionary contracting}\label{30007681}
\subsection{Definitions and mathematical statement}
Consider a market with one upstream incumbent, one possible upstream entrant and finitely many downstream retailers. Retailers purchase inputs under contracts specifying quantity-dependent payments and an exclusivity indicator. Let \(r\in\{0,1\}\) denote a regulatory regime: \(r=1\) permits exclusivity and nonlinear payments, while \(r=0\) prohibits exclusivity but preserves a specified feasible nonlinear-payment class. Let \(\theta\) collect consumer demand, production technologies, entrant fixed costs, investment opportunities and each firm's information. Firms bargain, choose investments and decide entry under a declared timing protocol. For each regime, let \(\mathcal E_r(\theta)\) contain equilibria consistent with that protocol; restrict comparisons to parameters for which both sets are nonempty. In equilibrium \(e\), define \(W(e)\) as consumer surplus plus producer profits, with profits measured net of real production, entry and investment expenditure. Determine welfare comparisons between regimes allowing retailers to rematch with suppliers, rather than mechanically preserving incumbent contracts. Separately define entry deterrence as lower entrant participation in regime \(1\) and productive coordination as lower real cost or higher investment-adjusted quality. Characterize when both effects coexist and which dominates welfare under a stated equilibrium-selection rule. The problem is to obtain tractable, estimable conditions within this model class and identify data on contract terms and switching that distinguish those mechanisms. It does not relabel previously solved exclusive-dealing examples as universal unresolved theorems.

\par\textbf{Formulation and scope.} The source verifies a broader question or research agenda; the population, estimand, equations and restrictions here are our scoped adaptation. This is not a quoted canonical conjecture, and the audit does not prove contemporary unsolved status for every special case. Independent math review tightened feasibility, existence, or estimand assumptions where required.

\par\textbf{Open-question provenance.} An explicit question or research agenda was verified at the source locator. This does not by itself certify that every later version remains unresolved \sref{30007681}{1}.

\par\textbf{Historical priority.} Earliest relevant explicit question or agenda passage located in this audit. No claim of first-ever formulation; earlier versions and later solutions were not exhaustively traced.

\subsection{Short English statement}
Compare exclusive-contract regimes in a specified supply-chain model, allowing coordination, entry deterrence and supplier rematching.

\subsection{Sources}
\begin{itemize}\raggedright
\item[\textnormal{[S1]}]\hypertarget{source:30007681:1}{} Robin S. Lee, Michael D. Whinston, Ali Yurukoglu. 2021. Structural Empirical Analysis of Contracting in Vertical Markets. Section 5, printed pp. 58–59, future research directions.
\url{https://robinlee.sites.fas.harvard.edu/papers/HandbookIO_Vol4_VerticalMarkets.pdf}.
DOI: 10.3386/w29282.
{\small\textit{Earliest verified question/agenda reference; historical priority qualified below. Richer contracts, partner choice and investment remain active modeling directions.}}
\end{itemize}
\end{document}
